% Plantilla creada por Sebastian Manrique.

% Las funciones \input representan la inserción de código .tex al documento principal, de esta manera podemos reducir varias líneas de código que en realidad, no se modifican con frecuencia.

\documentclass[letterpaper,11pt]{article}

\usepackage[spanish]{babel}
\usepackage[top=2cm,bottom=2cm,right=2cm,left=2cm]{geometry}
\usepackage{amssymb}
\usepackage{amsmath}
\usepackage{amsthm}
\usepackage{empheq}
\usepackage{mdframed}
\usepackage{booktabs}
\usepackage{lipsum}
\usepackage{graphicx}
\usepackage{color}
\usepackage{psfrag}
\usepackage{pgfplots}
\usepackage{bm}
\usepackage{multicol}
\usepackage{wrapfig}
\usepackage{float}
\usepackage{wrapfig}
\usepackage{multirow}
\usepackage{biblatex}
\usepackage{csquotes}
\usepackage{parskip}
\usepackage{cancel}
\usepackage{listing}
\usepackage{tikz}
\usepackage{verbatim}
\usepackage{subcaption}
\usepackage{enumitem}
\usepackage{xcolor}
\usepackage{xfrac}

\graphicspath{{fig/}}
\addbibresource{biblio.bib}
\usepgfplotslibrary{colorbrewer}
\pgfplotsset{width=8cm,compat=1.9}
\definecolor{negro}{RGB}{0,0,0}
\definecolor{azul}{RGB}{0,0,255}
\definecolor{verde}{RGB}{0,255,0}
\definecolor{cian}{RGB}{0,255,255}
\definecolor{rojo}{RGB}{255,0,0}
\definecolor{magenta}{RGB}{255,0,255}
\definecolor{amarillo}{RGB}{255,255,0}
\definecolor{blanco}{RGB}{255,255,255}

\definecolor{limegreen}{cmyk}{0.75,0,1,0}
\definecolor{orange}{cmyk}{0,0.8,0.95,0}
\definecolor{purple}{cmyk}{0.75,1,0,0}

\begin{document}

\begin{multicols}{2}

\section{Introducción}

$\Laplace{i}{\sin},\Fourier{i}{\dfrac{\sin}{\cos}}, \Laplace{}{\dfrac{\sin}{\cos}}, \Laplace{}{5}, \Laplace{i}{\sin}, \Laplace{}{kabfablfnsdojgbflsgen},$

\section{Sección 1} 

\begin{equation}
    Ecuaciones
\label{ecuacion}
\end{equation}

\begin{equation*} % Para ecuaciones sin número y que el igual se mantenga en la misma sangría
\begin{split}
    texto1 & = texto---2 \\
    texto---2 & = texto------3 \\
    texto------3 & = texto4
\end{split}
\end{equation*}

\section{Sección 2}

<<Texto>>

\end{multicols}

\subsection{Subsección 1}

<<Variación entre dos y una columna>> -------------------------------------------------------------------

\begin{multicols}{2}

\subsection{Subsección 2}

<<Texto>>

\subsubsection{Subsubsección 1}.

<<Texto>>

\begin{figure}[H]
    \centering
    \includegraphics[scale=0.32]{pollito.png}
    \caption{Inserción de una figura}
    \label{figura}
\end{figure}

\section{Conclusiones}

<<Conclusión>>

\end{multicols}

\section{Anexos}

<<Anexo>>

\end{document}